%! Unit 3: Polynomial Functions — Summative Assessment — Unit Test (blank)
\documentclass[10pt]{article}
\usepackage{algebra2-article}
\usepackage{algebra2-boxes}
% Coordinate grid: {xmin}{xmax}{ymin}{ymax}
\newcommand{\lgrid}[4]{%
  \draw[step=1, linegray!40, very thin] (#1,#3) grid (#2,#4);
  \draw[->, semithick, charcoal] (#1,0)--(#2,0) node[below right, font=\scriptsize]{$x$};
  \draw[->, semithick, charcoal] (0,#3)--(0,#4) node[left, font=\scriptsize]{$y$};}
\newcommand{\dpt}[2]{\fill[charcoal] (#1,#2) circle (3pt);}

% Part-divider strip (headlinebox takes one arg: the background color)
\newcommand{\parthead}[1]{%
  \vspace{6pt}\noindent
  \begin{headlinebox}{forest}\color{white}\bfseries #1\end{headlinebox}%
  \vspace{4pt}}

\begin{document}

\pageheader{Unit 3: Polynomial Functions}{Unit Test}
\namedateperiod

\vspace{0.06in}
\noindent\textbf{Instructions:} Show all work. No notes unless your teacher says so.

% ======================================================================
\parthead{Part A --- Vocabulary (8 pts)}
Write the letter of the best description next to each term.

\vspace{4pt}
\noindent
\begin{minipage}[t]{0.40\textwidth}
\begin{enumerate}[leftmargin=2.2em, itemsep=7pt, label=\arabic*.]
  \item Multiplicity of a zero \blank{1.2cm}
  \item End behavior \blank{1.2cm}
  \item Rational Root Theorem \blank{1.2cm}
  \item Degree \blank{1.2cm}
  \item Turning point \blank{1.2cm}
  \item Fundamental Theorem of Algebra \blank{1.2cm}
  \item Leading coefficient \blank{1.2cm}
  \item Factor Theorem \blank{1.2cm}
\end{enumerate}
\end{minipage}%
\hfill
\begin{minipage}[t]{0.57\textwidth}
\small
\begin{description}[leftmargin=1.4em, itemsep=3pt]
  \item[(A)] The number multiplying the highest-power term; its sign decides the right-hand arm.
  \item[(B)] The greatest exponent on the variable once the polynomial is in standard form.
  \item[(C)] How many times the factor $(x-r)$ appears --- odd means the graph crosses at $r$, even means it touches.
  \item[(D)] A point where the graph changes from increasing to decreasing, or the reverse.
  \item[(E)] What the graph does at its far left and far right, as $x\to\pm\infty$.
  \item[(F)] $(x-r)$ is a factor of $f(x)$ exactly when $f(r)=0$.
  \item[(G)] Every rational zero of a polynomial with integer coefficients looks like $\pm\frac{p}{q}$, where $p$ divides the constant term and $q$ divides the leading coefficient.
  \item[(H)] A polynomial of degree $n\ge 1$ has exactly $n$ complex zeros, counted with multiplicity.
\end{description}
\end{minipage}

% ======================================================================
\parthead{Part B --- Multiple Choice (12 pts)}
Circle the best answer.

\begin{enumerate}[leftmargin=1.9em, itemsep=9pt]
  \item The end behavior of $f(x)=2x^3-5x+1$ is
    \hfill (A) both arms up \quad (B) both arms down \quad (C) left down, right up \quad (D) left up, right down
  \item $(3x+4)^2=$
    \hfill (A) $9x^2+16$ \quad (B) $9x^2+24x+16$ \quad (C) $9x^2+12x+16$ \quad (D) $9x^2-16$
  \item Factored completely, $x^3-8$ is
    \hfill (A) $(x-2)(x^2+2x+4)$ \quad (B) $(x-2)(x^2-2x+4)$ \quad (C) $(x-2)^3$ \quad (D) prime
  \item The remainder when $f(x)=x^3+2x^2-5$ is divided by $x+1$ is
    \hfill (A) $-4$ \quad (B) $-2$ \quad (C) $4$ \quad (D) $0$
  \item Which value is \emph{not} a possible rational zero of $f(x)=2x^3-x^2+5x-3$?
    \hfill (A) $-3$ \quad (B) $\frac{3}{2}$ \quad (C) $\frac{2}{3}$ \quad (D) $1$
  \item A polynomial with \emph{real} coefficients has degree $4$ and zeros $-3$, $1$, and $2i$. Its fourth zero must be
    \hfill (A) $-2i$ \quad (B) $2i$ \quad (C) $-2$ \quad (D) $\frac12$
\end{enumerate}

% ======================================================================
\parthead{Part C --- Short Answer \& Computation (40 pts)}
Show your work.

\begin{enumerate}[leftmargin=1.9em, itemsep=10pt]
  \item \textbf{Read the graph.} The polynomial function $f$ is graphed below. Each gridline is one unit.
        \begin{center}
        \begin{tikzpicture}[scale=0.5]
          \lgrid{-4}{3}{-5}{5}
          \begin{scope}
            \clip (-4,-5) rectangle (3,5);
            \draw[very thick, forest, domain=-3.2:1.5, samples=180, smooth]
              plot (\x,{(\x)*(\x)*(\x) + 3*(\x)*(\x) - 4});
          \end{scope}
          \foreach \x in {-3,-2,-1,1,2}
            \draw[charcoal] (\x,-0.15)--(\x,0.15) node[below=1pt, font=\tiny]{$\x$};
          \dpt{-2}{0}
          \dpt{1}{0}
          \dpt{0}{-4}
        \end{tikzpicture}
        \end{center}
        (a) Degree: even or odd? \blank{1.4cm} \quad Leading coefficient: positive or negative? \blank{1.8cm}
        \\[4pt]
        (b) As $x\to-\infty$, $f(x)\to$ \blank{1.4cm}; \quad as $x\to+\infty$, $f(x)\to$ \blank{1.4cm}
        \\[4pt]
        (c) Zeros, each with its multiplicity, and whether the graph crosses or touches there:
        \blank{7.0cm}
        \\[4pt]
        (d) $y$-intercept: \blank{1.8cm}; \quad number of turning points: \blank{1.2cm}
        \\[4pt]
        (e) Relative maximum: \blank{2.0cm}; \quad relative minimum: \blank{2.0cm};
        \\[4pt]
        \phantom{(e)} absolute maximum or minimum: \blank{4.0cm}
        \\[4pt]
        (f) Interval(s) where $f$ is decreasing: \blank{2.6cm}; \quad domain: \blank{2.0cm}; \quad range: \blank{2.0cm}

  \item \textbf{Operations.} Simplify. Write each answer in standard form.
        \\[4pt]
        (a) $(5x^3-2x^2+7)-(2x^3+x^2-3x)$
        \vspace{1.5cm}
        \\
        (b) $(3x-2)(x^2-5x+4)$
        \vspace{1.6cm}
        \\
        (c) For your answer to (b), give the degree \blank{1.0cm}, the leading coefficient \blank{1.0cm},
        and the end behavior \blank{4.2cm}

  \item \textbf{Factor completely.} Every factor must be prime over the integers.
        \\[4pt]
        (a) $3x^4-48x^2$ \hspace{1.5cm} (b) $x^3+125$ \hspace{1.5cm} (c) $x^3-2x^2-9x+18$
        \vspace{2.8cm}

  \item \textbf{Division and the Remainder Theorem.}
        \\[4pt]
        (a) Divide using \emph{long division} and write the answer as $Q(x)+\dfrac{R}{D(x)}$:
        \[ (x^3-3x^2+5)\div(x+1) \]
        \vspace{2.6cm}
        \\
        (b) Use the \emph{Remainder Theorem} to find the remainder when $g(x)=x^3-2x^2-4x+8$ is divided
        by $x-2$. Is $x-2$ a factor of $g$? Explain in one sentence.
        \vspace{2.0cm}

  \item \textbf{Synthetic division.} You are told that $x=-1$ is a zero of
        \[ f(x)=x^3-x^2-10x-8. \]
        Use synthetic division to divide out that zero, factor $f$ completely, and list \emph{all} of
        its zeros.
        \vspace{3.0cm}

  \item \textbf{Rational Root Theorem.} For $f(x)=3x^3-x^2-12x+4$:
        \\[4pt]
        (a) List every possible rational zero.
        \vspace{1.5cm}
        \\
        (b) Find all real zeros of $f$, showing the test and the division you used.
        \vspace{3.0cm}

  \item \textbf{Solve over the complex numbers.} $x^3-2x^2+9x-18=0$
        \\[4pt]
        (a) Before solving: how many solutions must this equation have, counted with multiplicity, and
        which theorem says so? \blank{6.0cm}
        \\[4pt]
        (b) Solve. Give every solution.
        \vspace{2.6cm}
        \\
        (c) How many of your solutions are real, and how many $x$-intercepts does the graph of
        $y=x^3-2x^2+9x-18$ therefore have? \blank{5.0cm}

  \item \textbf{Build a polynomial from its zeros.} Write a polynomial function of \emph{least degree}
        with \emph{real} coefficients whose zeros are $3$ and $-i$. Give your answer in factored form
        \emph{and} in standard form.
        \vspace{2.8cm}
\end{enumerate}

% ======================================================================
\parthead{Part D --- Extended Response (12 pts)}
Explain your reasoning in full sentences.

\begin{enumerate}[leftmargin=1.9em, itemsep=16pt]
  \item \textbf{The graph plan.} Consider $f(x)=3(x+3)(x-1)^2(x-2)$. Do \emph{not} expand it.
        \textbf{(a)} State the degree and the leading coefficient, and explain how you got each one
        from the factored form. What end behavior do they force?
        \textbf{(b)} List the zeros with their multiplicities, and say whether the graph crosses or
        touches the $x$-axis at each.
        \textbf{(c)} Find the $y$-intercept.
        \textbf{(d)} Build a sign chart and state every interval on which $f(x)>0$.
        \textbf{(e)} At most how many turning points can this graph have, and why?
        \vspace{6.0cm}

  \item \textbf{Modeling.} An open-top box is made from a $14$ in by $8$ in sheet of cardboard by
        cutting a square of side $x$ from each corner and folding up the sides.
        \textbf{(a)} Write the volume $V(x)$ in factored form, then expand it into standard form.
        \textbf{(b)} State the feasible domain of the model and justify it from the sheet's dimensions.
        \textbf{(c)} Find $V(1)$ and state its units.
        \textbf{(d)} Check that $x=7$ makes $V(x)=0$. Explain why $7$ is a genuine zero of the
        polynomial but not a possible cut length.
        \vspace{5.5cm}
\end{enumerate}

\end{document}
